% GENERATED FILE. Edit canonical lessons, then run npm run generate:books.
% canonical-source-hash: fnv1a64:6bcbbc294d006003
% canonical-lessons: MW-C269-bahano, MW-C269-inkar, MW-C269-arno, MW-C269-jid, MW-C269-tokno

\chapter{Not Now}
\label{ch:mw-not-now}
\hlchaptermodality{269}

\begin{chapteropening}
\textbf{By the end of this chapter:} \emph{I can turn something down and refuse.}

Complete the last lesson of chapter 269: I can turn something down and refuse.
\end{chapteropening}

\section[\texorpdfstring{\mw{बहानो} (bahāno)}{बहानो (bahāno)}]{\mw{बहानो} (bahāno) — an excuse}

\label{lesson:MW-C269-bahano}



\begin{quote}
Before the new one: say the Marwadi for a meaning, then the Marwadi for a good word for someone.
\end{quote}

\subsection*{You'll want to know: \mw{बहानो}}
\textbf{\mw{बहानो}} — \emph{bahāno} — \textquotedblleft{}an excuse\textquotedblright{}.

\begin{quote}
\mw{यो} \mw{बहानो} \mw{साचो} \mw{कोनी।}
\end{quote}

\begin{quote}
\emph{yo bahāno sācho konī.}
\end{quote}

\begin{quote}
This excuse isn't true.
\end{quote}

A masculine noun.

\begin{cousinweb}
\textbf{\mw{बहानो}} (\emph{bahāno}) is usually traced to Persian \emph{bahāna}, \textquotedblleft{}a pretext\textquotedblright{}; Marwadi gives it the -\mw{ो} ending of a masculine noun.
\end{cousinweb}

\begin{grammarlens}[title={What you've built}]
One more word for asking and answering.
\end{grammarlens}

\subsection*{Your turn}
\begin{itemize}
\raggedright
  \item \emph{Say it:} \emph{bahāno}
  \item \emph{Say it:} \emph{bahāno}, once more
  \item \emph{Say it:} \emph{yo bahāno sācho konī.} — This excuse isn't true.
  \item \emph{Recall:} say the Marwadi for a meaning, then the Marwadi for a good word for someone, then say \emph{bahāno} again
\end{itemize}

\begin{tcolorbox}[breakable,colback=teal!4,colframe=teal!35!black,title={Before you move on}]
What does \mw{बहानो} mean? (\textquotedblleft{}An excuse\textquotedblright{}.) Say it once more.
\end{tcolorbox}
\section[\texorpdfstring{\mw{इनकार} (inkār)}{इनकार (inkār)}]{\mw{इनकार} (inkār) — a refusal}

\label{lesson:MW-C269-inkar}



\begin{quote}
Before the new one: say the Marwadi for a good word for someone, then the Marwadi for an excuse.
\end{quote}

\subsection*{You'll want to know: \mw{इनकार}}
\textbf{\mw{इनकार}} — \emph{inkār} — \textquotedblleft{}a refusal\textquotedblright{}.

\begin{quote}
\mw{म्हैं} \mw{इनकार} \mw{कियो।}
\end{quote}

\begin{quote}
\emph{mhaĩ inkār kiyo.}
\end{quote}

\begin{quote}
I refused.
\end{quote}

A masculine noun.

\begin{cousinweb}
\textbf{\mw{इनकार}} (\emph{inkār}) is usually traced, through Persian \emph{inkār}, to Arabic \emph{inkār}, \textquotedblleft{}denial\textquotedblright{}. With \textbf{\mw{कियो}} it says you refused: \textbf{\mw{म्हैं} \mw{इनकार} \mw{कियो}}.
\end{cousinweb}

\begin{grammarlens}[title={What you've built}]
One more word for asking and answering.
\end{grammarlens}

\subsection*{Your turn}
\begin{itemize}
\raggedright
  \item \emph{Say it:} \emph{inkār}
  \item \emph{Say it:} \emph{inkār}, once more
  \item \emph{Say it:} \emph{mhaĩ inkār kiyo.} — I refused.
  \item \emph{Recall:} say the Marwadi for a good word for someone, then the Marwadi for an excuse, then say \emph{inkār} again
\end{itemize}

\begin{tcolorbox}[breakable,colback=teal!4,colframe=teal!35!black,title={Before you move on}]
What does \mw{इनकार} mean? (\textquotedblleft{}A refusal\textquotedblright{}.) Say it once more.
\end{tcolorbox}
\section[\texorpdfstring{\mw{अड़णो} (aṛṇo)}{अड़णो (aṛṇo)}]{\mw{अड़णो} (aṛṇo) — to dig in one's heels; to insist}

\label{lesson:MW-C269-arno}



\begin{quote}
Before the new one: say the Marwadi for an excuse, then the Marwadi for a refusal.
\end{quote}

\subsection*{You'll want to know: \mw{अड़णो}}
\textbf{\mw{अड़णो}} — \emph{aṛṇo} — \textquotedblleft{}to dig in one's heels; to insist\textquotedblright{}.

\begin{quote}
\mw{छोरो} \mw{अड़ै} \mw{है} \mw{अर} \mw{स्कूल} \mw{कोनी} \mw{जावै।}
\end{quote}

\begin{quote}
\emph{chhoro aṛai hai ar skūl konī jāvai.}
\end{quote}

\begin{quote}
The boy digs in his heels and won't go to school.
\end{quote}

Past (he …): \mw{वो} \mw{अड़्यो}, \emph{vo aṛyo}.

\begin{cousinweb}
\textbf{\mw{अड़णो}} (\emph{aṛṇo}) is inherited from a reconstructed \emph{*aḍ}, \textquotedblleft{}obstruct, stop\textquotedblright{}; CDIAL gives no Marwari form in this entry, but lists Hindi \emph{aṛnā}, \textquotedblleft{}to come to a stop\textquotedblright{}, and Kumaoni \emph{aṛṇo}, \textquotedblleft{}to persist\textquotedblright{} (CDIAL 187). Someone who stops and will not move on is digging in.
\end{cousinweb}

\begin{grammarlens}[title={What you've built}]
One more word for asking and answering.
\end{grammarlens}

\subsection*{Your turn}
\begin{itemize}
\raggedright
  \item \emph{Say it:} \emph{aṛṇo}
  \item \emph{Say it:} \emph{aṛṇo}, once more
  \item \emph{Say it:} \emph{chhoro aṛai hai ar skūl konī jāvai.} — The boy digs in his heels and won't go to school.
  \item \emph{Recall:} say the Marwadi for an excuse, then the Marwadi for a refusal, then say \emph{aṛṇo} again
\end{itemize}

\begin{tcolorbox}[breakable,colback=teal!4,colframe=teal!35!black,title={Before you move on}]
What does \mw{अड़णो} mean? (\textquotedblleft{}To dig in one's heels; to insist\textquotedblright{}.) Say it once more.
\end{tcolorbox}
\section[\texorpdfstring{\mw{जिद} (jid)}{जिद (jid)}]{\mw{जिद} (jid) — stubbornness}

\label{lesson:MW-C269-jid}



\begin{quote}
Before the new one: say the Marwadi for a refusal, then the Marwadi for to dig in one's heels.
\end{quote}

\subsection*{You'll want to know: \mw{जिद}}
\textbf{\mw{जिद}} — \emph{jid} — \textquotedblleft{}stubbornness\textquotedblright{}.

\begin{quote}
\mw{जिद} \mw{मत} \mw{करो।}
\end{quote}

\begin{quote}
\emph{jid mat karo.}
\end{quote}

\begin{quote}
Don't be stubborn.
\end{quote}

A feminine noun.

\begin{cousinweb}
\textbf{\mw{जिद}} (\emph{jid}) is usually traced, through Persian \emph{zidd}, to Arabic \emph{ḍidd}, \textquotedblleft{}the opposite, opposition\textquotedblright{}.
\end{cousinweb}

\begin{grammarlens}[title={What you've built}]
One more word for asking and answering.
\end{grammarlens}

\subsection*{Your turn}
\begin{itemize}
\raggedright
  \item \emph{Say it:} \emph{jid}
  \item \emph{Say it:} \emph{jid}, once more
  \item \emph{Say it:} \emph{jid mat karo.} — Don't be stubborn.
  \item \emph{Recall:} say the Marwadi for a refusal, then the Marwadi for to dig in one's heels, then say \emph{jid} again
\end{itemize}

\begin{tcolorbox}[breakable,colback=teal!4,colframe=teal!35!black,title={Before you move on}]
What does \mw{जिद} mean? (\textquotedblleft{}Stubbornness\textquotedblright{}.) Say it once more.
\end{tcolorbox}
\section[\texorpdfstring{\mw{टोकणो} (ṭokṇo)}{टोकणो (ṭokṇo)}]{\mw{टोकणो} (ṭokṇo) — to interrupt; to object}

\label{lesson:MW-C269-tokno}



\begin{quote}
Before the new one: say the Marwadi for to dig in one's heels, then the Marwadi for stubbornness.
\end{quote}

\subsection*{You'll want to know: \mw{टोकणो}}
\textbf{\mw{टोकणो}} — \emph{ṭokṇo} — \textquotedblleft{}to interrupt; to object\textquotedblright{}.

\begin{quote}
\mw{मास्टर} \mw{बोलै} \mw{है}, \mw{मत} \mw{टोको।}
\end{quote}

\begin{quote}
\emph{māsṭar bolai hai, mat ṭoko.}
\end{quote}

\begin{quote}
The teacher is speaking: don't interrupt.
\end{quote}

Past (I …): \mw{म्हैं} \mw{थारी} \mw{बात} \mw{कोनी} \mw{टोकी}, \emph{mhaĩ thārī bāt konī ṭokī}.

\begin{cousinweb}
\textbf{\mw{टोकणो}} (\emph{ṭokṇo}) is inherited from a reconstructed \emph{*ṭōkk}, \textquotedblleft{}strike, bite\textquotedblright{}; CDIAL gives no Marwari form in this entry, but lists Hindi \emph{ṭoknā}, \textquotedblleft{}to check, blame, question\textquotedblright{} (CDIAL 5476).
\end{cousinweb}

\begin{grammarlens}[title={What you've built}]
One more word for asking and answering.
\end{grammarlens}

\subsection*{Your turn}
\begin{itemize}
\raggedright
  \item \emph{Say it:} \emph{ṭokṇo}
  \item \emph{Say it:} \emph{ṭokṇo}, once more
  \item \emph{Say it:} \emph{māsṭar bolai hai, mat ṭoko.} — The teacher is speaking: don't interrupt.
  \item \emph{Recall:} say the Marwadi for to dig in one's heels, then the Marwadi for stubbornness, then say \emph{ṭokṇo} again
\end{itemize}

\begin{tcolorbox}[breakable,colback=teal!4,colframe=teal!35!black,title={Before you move on}]
What does \mw{टोकणो} mean? (\textquotedblleft{}To interrupt; to object\textquotedblright{}.) Say it once more.
\end{tcolorbox}
